\chapter{Cyclic Torsor Theories and Witt Vectors}
\label{appendix:witt_vector_algebra}
\label{appendix:cyclic_torsor_theories}

This appendix contains the supporting theory and proofs for the three cyclic
torsor descriptions stated in Chapter~2. We begin with the common
group-scheme construction and the classical Kummer and Artin--Schreier cases.
We then develop the finite-length \(p\)-typical Witt-vector algebra needed for
Artin--Schreier--Witt theory, followed by its torsor classification and
normal-form arguments. The final section isolates the weighted-homogeneity
property used in the cyclic symmetric-power calculation. The Schmid--Witt
ramification-break formula is stated, referenced, and applied in Chapter~2;
we do not reproduce its proof.

\providecommand{\bW}{\mathbb{W}}

\section{The common group-scheme construction}
\label{appendix:common_group_scheme_construction}

The three cyclic theories used in Chapter~2 are instances of the following
construction.

\subsection{The kernel torsor construction}

Let $H$ be a commutative group scheme over a scheme $S$, written additively,
and let $\varphi:H\to H$ be a homomorphism. Its scheme-theoretic kernel
$K:=\ker(\varphi)$ is defined by the Cartesian square
\[
\begin{tikzcd}
K \arrow[r] \arrow[d] & H \arrow[d,"\varphi"] \\
S \arrow[r,"0_H"'] & H.
\end{tikzcd}
\]
Thus, for every $S$-scheme $U$,
\begin{equation}\label{eq:kernel_on_U_valued_points}
K(U)=\ker(\varphi_U)
    =\{k\in H(U):\varphi_U(k)=0_{H(U)}\},
\end{equation}
where $H(U)=\Hom_S(U,H)$ and $\varphi_U$ is induced by $\varphi$.

For an $S$-scheme $X$ and a section $a\in H(X)$, define the fiber
\begin{equation}\label{eq:fiber_torsor_definition}
T_a:=X\times_{a,H,\varphi}H,
\end{equation}
equivalently by the Cartesian square
\[
\begin{tikzcd}
T_a \arrow[r] \arrow[d] & H \arrow[d,"\varphi"] \\
X \arrow[r,"a"'] & H.
\end{tikzcd}
\]
For every $X$-scheme $U\to X$, with $a_U$ the pullback of $a$,
\[
T_a(U)=\{y\in H(U):\varphi_U(y)=a_U\}.
\]
Thus ``the torsor $\varphi(Y)=a$'' means the scheme of solutions after every
base change to $X$; it has no preferred identity element in general.

Write $H_X:=H\times_S X$, $K_X:=K\times_S X$, let
$\varphi_X:H_X\to H_X$ be the base change of $\varphi$, and let
$\iota_K:K_X\to H_X$ and $\iota_a:T_a\to H_X$ be the canonical maps.
Addition on $H_X$ restricts to the action
\begin{equation}\label{eq:kernel_action_on_fiber}
K_X\times_X T_a\longrightarrow T_a,
\qquad (k,y)\longmapsto y+k.
\end{equation}
It is the top arrow in the commutative diagram
\[
\begin{tikzcd}[column sep=large, row sep=large]
K_X\times_X T_a
    \arrow[r,"{(k,y)\mapsto y+k}"]
    \arrow[d,"{\iota_K\times_X\iota_a}"']
  & T_a \arrow[d,"\iota_a"] \\
H_X\times_X H_X
    \arrow[r,"+_X"']
  & H_X.
\end{tikzcd}
\]
Indeed,
\[
\varphi_U(y+k)
=\varphi_U(y)+\varphi_U(k)
=a_U+0_{H(U)}
=a_U,
\]
so translation by $K$ preserves the fiber $T_a$. The action is simply
transitive: the morphism
\begin{equation}\label{eq:torsor_difference_isomorphism}
K_X\times_X T_a\longrightarrow T_a\times_X T_a,
\qquad (k,y)\longmapsto (y,y+k),
\end{equation}
is an isomorphism with inverse
\[
(y_1,y_2)\longmapsto (y_2-y_1,y_1).
\]
This inverse is well-defined because
\[
\varphi_U(y_2-y_1)
=\varphi_U(y_2)-\varphi_U(y_1)
=a_U-a_U
=0_{H(U)}.
\]
If $\varphi$ is finite \'etale and surjective, then its base change
$T_a\to X$ is finite \'etale and surjective. Hence $T_a$ is a $K_X$-torsor,
namely
\[
T_a=a^*(H\xrightarrow{\varphi}H)
\]
as a pullback of the $K$-torsor $\varphi:H\to H$.

Suppose that $a'=a+\varphi_X(b)$ for some $b\in H(X)$. Translation by $b$
defines
\[
\tau_b:T_a\longrightarrow T_{a'},
\qquad y\longmapsto y+b,
\]
because
\[
\varphi_U(y+b_U)
=a_U+\varphi_U(b_U)
=a'_U,
\]
and it is $K_X$-equivariant since
$\tau_b(y+k)=y+k+b=\tau_b(y)+k$. Its inverse is $\tau_{-b}$, so $T_a$ depends
only on the class of $a$ modulo $\varphi_X(H(X))$.

When $\varphi$ is finite \'etale and surjective, the exact sequence on
$X_{\mathrm{\acute et}}$
\[
0\longrightarrow K_X\longrightarrow H_X
\xrightarrow{\varphi_X}H_X\longrightarrow0
\]
has connecting map $a\mapsto[T_a]$ and induces the injection
\begin{equation}\label{eq:common_torsor_connecting_map}
H(X)/\varphi_X(H(X))\lhook\joinrel\longrightarrow
H^1_{\mathrm{\acute et}}(X,K_X),
\qquad [a]\longmapsto[T_a].
\end{equation}
The class $[T_a]$ is trivial exactly when $a\in\varphi_X(H(X))$, equivalently
when $\varphi(Y)=a$ has a global solution.

For a multiplicative group scheme the same construction is written
multiplicatively: $K(U)=\{k\in H(U):\varphi_U(k)=1_{H(U)}\}$, the action is
$(k,y)\mapsto ky$, and replacing $a$ by $a\varphi_X(b)$ gives the isomorphic
torsor $y\mapsto by$.

\subsection{The three cyclic theories}

For an arbitrary $S$-scheme $X$, the fiber over a section $a\in H(X)$ is
\[
T_a=\varphi_X^{-1}(a)
   =X\times_{a,H,\varphi}H.
\]
For nonaffine $X$, explicit equations are interpreted using the structure
sheaf $\Oo_X$. In all three applications in Chapter~2, however, the base is
affine: first $\Spec k((x))$ and then
$\bG_m=\Spec k[x,x^{-1}]$. We therefore work affinely from now on and write
\[
X=\Spec A,
\]
where the capital letter $A$ is the coordinate ring of $X$, while the
lowercase letter $a$ is the section specifying the fiber. More generally, if
$S=\Spec R$ and $H=\Spec C$, then
\[
T_a=\Spec\bigl(A\otimes_C C\bigr),
\]
where $C$ acts on $A$ through $a^\#:C\to A$ and on the second copy of $C$
through $\varphi^\#:C\to C$.

\paragraph{Kummer.}
Take $H=\bG_m$ and $\varphi(z)=z^e$. A section is an element
$a\in\bG_m(X)=A^\times$, and its fiber is
\[
T_a=\Spec\bigl(A[Y]/(Y^e-a)\bigr),
\]
where $Y$ is automatically a unit because $a$ is a unit. The kernel is
\[
\mu_{e,X}=\Spec\bigl(A[t,t^{-1}]/(t^e-1)\bigr),
\]
acting by
$(\zeta,Y)\mapsto\zeta Y$. If $e$ is invertible on $S$ and a primitive root
$\zeta_e\in\mu_e(S)$ is chosen, then
$\overline j\mapsto\zeta_e^j$ identifies $(\Z/e\Z)_S$ with $\mu_e$; only
this identification, not the $\mu_e$-torsor, depends on the choice.

\paragraph{Artin--Schreier.}
Let $A$ be an $\F_p$-algebra, take $H=\bG_a$, and put
$\varphi(Y)=Y^p-Y$. A section is an element $a\in\bG_a(X)=A$, and its
fiber is
\[
T_a=\Spec\bigl(A[Y]/(Y^p-Y-a)\bigr).
\]
The kernel is the constant group scheme $(\F_p)_X$, acting by
$(c,Y)\mapsto Y+c$.

\paragraph{Artin--Schreier--Witt.}
Let $A$ be an $\F_p$-algebra. The third theory takes $H=\bW_r$ and
\[
\varphi=\wp:=F-\mathrm{id},
\qquad
\wp\vec Y=F\vec Y\ominus\vec Y.
\]
For a section $\vec a\in W_r(A)$, the fiber is explicitly
\[
T_{\vec a}
=\Spec\!\left(
  A[Y_0,\dots,Y_{r-1}]
  \big/\bigl(\wp\vec Y\ominus\vec a\bigr)
\right),
\qquad \wp\vec Y=\vec a.
\]
Here the parenthesized Witt vector denotes the ideal generated by its $r$
coordinate polynomials. Its kernel is
$W_r(\F_p)_X\cong(\Z/p^r\Z)_X$, acting by Witt translation
$(\vec c,\vec Y)\mapsto\vec Y\oplus\vec c$.

This paragraph states the affine model in advance for comparison with the
two preceding theories. The operations $F$, $\oplus$, and $\ominus$, and the
coordinate polynomials occurring in the displayed quotient, are defined in
\Cref{appendix:witt_vector_algebra_basics}. The fact that this fiber is
finite \'etale and that the displayed kernel action makes it a torsor is
proved in \Cref{lemma:asw_etale_torsor}.

\section{Kummer and Artin--Schreier theory}
\label{appendix:kummer_and_as_proofs}

Let $K$ be a complete discretely valued field with residue characteristic
$p>0$.

\subsection{Kummer theory}

If $(e,p)=1$ and $K$ contains $\mu_e$, the Kummer sequence
\[
1\longrightarrow\mu_e\longrightarrow\bG_m
\xrightarrow{(-)^e}\bG_m\longrightarrow1
\]
and Hilbert~90 give
\[
H^1_{\mathrm{\acute et}}(\Spec K,\mu_e)
\cong K^\times/(K^\times)^e.
\]
Equivalently, since the base is a field, this \'etale cohomology group is the
continuous Galois cohomology group
\[
H^1_{\mathrm{\acute et}}(\Spec K,\mu_e)
\cong
H^1_{\mathrm{cont}}\!\left(\Gal(K^{\mathrm{sep}}/K),
\mu_e(K^{\mathrm{sep}})\right).
\]
The class of $a\in K^\times$ is represented by $T^e=a$, with
$\zeta\in\mu_e$ acting by $T\mapsto\zeta T$.

\begin{proof}[Proof of \Cref{theorem:kummer_ramification}]
Write $a=\pi^m u$ with $u\in\Oo_K^\times$. Replacing $a$ by an element of
the same class in $K^\times/(K^\times)^n$ changes $m$ by a multiple of $n$.
If $n\mid m$, we may therefore assume that $a=u$ is a unit. Since $n$ is
invertible in the residue field, Hensel's lemma shows that $u$ is an $n$-th
power in $K$ precisely when its residue is an $n$-th power; this is the
trivial case. Otherwise the extension is obtained by lifting the corresponding
separable residue-field extension and is unramified.

In general, if $y^n=a$, then $v_L(y)=m/n$ with respect to the extension of
$v_K$. Hence the subgroup of the value group generated by $\Z$ and $m/n$ has
index $n/\gcd(n,m)$ over $\Z$, which is the ramification index. When
$\gcd(n,m)=1$ this index is $n$, so the extension is totally ramified; it is
tame because $(n,p)=1$. The converse is the Kummer classification supplied by
the displayed cohomological isomorphism. See
\cite[Proposition~1.83]{koch1997algebraic} for the general statement.
\end{proof}

\subsection{Artin--Schreier theory}

In characteristic $p$, the Artin--Schreier sequence
\[
0\longrightarrow\F_p\longrightarrow\bG_a
\xrightarrow{\wp}\bG_a\longrightarrow0,
\qquad \wp(z)=z^p-z,
\]
and additive Hilbert~90 give
\[
H^1(K,\Z/p\Z)\cong K/\wp(K).
\]
The class of $a\in K$ is represented by $X^p-X=a$, with
$c\in\F_p$ acting by $X\mapsto X+c$.

\begin{proof}[Proof of \Cref{theorem:artin_schreier_ramification}]
If $v_K(a)>0$, the convergent series
$u=-\sum_{i\geq0}a^{p^i}$ satisfies $u^p-u=a$, so the class is trivial. If
$v_K(a)=0$, the same argument removes every term of positive valuation. The
extension is therefore determined by the Artin--Schreier class of the residue
of $a$: it is trivial when this residue class vanishes and otherwise is the
unramified cyclic extension of degree $p$.

Suppose now that $v_K(a)=-m<0$ and $p\nmid m$, and let $y^p-y=a$. Since the
residue class cannot be removed by an Artin--Schreier change of variable, the
extension has degree $p$. Normalize $v_L$ so that $v_L|_K=pv_K$. The equation
gives $v_L(y)=-m$, so the extension is totally ramified. Choose integers
$1\leq c\leq p-1$ and $b$ with $-cm+bp=1$. Then
$\varpi_L=y^c\pi^b$ is a uniformizer. For a nontrivial
$\sigma\in\Gal(L/K)$, after rescaling the generator we have
$\sigma(y)=y+1$, and
\[
v_L\bigl(\sigma(\varpi_L)-\varpi_L\bigr)
=v_L\bigl(\pi^b((y+1)^c-y^c)\bigr)
=bp-(c-1)m=m+1.
\]
Thus the unique lower ramification break is $m$; for a cyclic group of order
$p$ this is also the unique upper break. Finally, the Artin--Schreier
cohomological classification gives every cyclic extension of degree $p$.
The reduced polynomial representative used in the last assertion is
constructed in the proof of \Cref{lemma:realization_over_Gm}(2).
\end{proof}

\begin{lemma}[Classical cyclic equations are \'etale torsors]
\label{lemma:classical_cyclic_equations_etale_torsors}
Let $A$ be a ring.
\begin{enumerate}[label=(\roman*)]
    \item If $n$ is invertible in $A$ and $a\in A^\times$, then
    \[
    A[T]/(T^n-a)
    \]
    is finite free over $A$, with basis $1,T,\dots,T^{n-1}$, and is
    \'etale over $A$. Multiplication of $T$ makes its spectrum a
    $\mu_n$-torsor over $\Spec A$.
    \item If $A$ is an $\F_p$-algebra and $a\in A$, then
    \[
    A[X]/(X^p-X-a)
    \]
    is finite free over $A$, with basis $1,X,\dots,X^{p-1}$, and is
    \'etale over $A$. Translation of $X$ makes its spectrum a
    $\Z/p\Z$-torsor over $\Spec A$.
\end{enumerate}
\end{lemma}

\begin{proof}
Both defining polynomials are monic, giving the stated bases and finite
freeness. In the Kummer algebra, $T$ is a unit because $T^n=a\in A^\times$;
therefore the derivative $nT^{n-1}$ is a unit. The derivative of
$X^p-X-a$ is $-1$. Thus both algebras are \'etale over $A$. They are the
pullbacks from Appendix~A.1 of $(-)^n:\bG_m\to\bG_m$ and
$\wp:\bG_a\to\bG_a$, respectively, so translation by the corresponding
kernels $\mu_n$ and $\F_p\cong\Z/p\Z$ gives the asserted torsor structures.
\end{proof}

\section{Witt-vector algebra}
\label{appendix:witt_vector_algebra_basics}

\subsection{Ring structure and ghost coordinates}

For $n \geq 0$, the $n$-th \textbf{Witt polynomial} in the variables $X_0, \dots, X_n$
is
\[
w_n(X_0, \dots, X_n) \;=\; \sum_{j=0}^{n} p^j X_j^{p^{n-j}}
\;=\; X_0^{p^n} + p X_1^{p^{n-1}} + \dots + p^n X_n .
\]

\begin{theorem}[Witt; {\cite[II, \S6, Theorem~6]{serreLF}}, {\cite[\S1]{Rabinoff2014Witt}}]
\label{theorem:witt_ring_structure}
There exists a unique family of polynomials
$S_n, P_n \in \Z[X_0, \dots, X_n; Y_0, \dots, Y_n]$, $n \geq 0$, such that
\[
\begin{aligned}
w_n(S_0, \dots, S_n)
    &= w_n(X_0,\dots,X_n) + w_n(Y_0,\dots,Y_n),\\
w_n(P_0, \dots, P_n)
    &= w_n(X_0,\dots,X_n) \cdot w_n(Y_0,\dots,Y_n).
\end{aligned}
\]
for all $n$. For any commutative ring $A$, the operations
\[
\vec a \oplus \vec b := \bigl(S_n(\vec a; \vec b)\bigr)_{0 \le n \le r-1},
\qquad
\vec a \odot \vec b := \bigl(P_n(\vec a; \vec b)\bigr)_{0 \le n \le r-1}
\]
make the set $A^r$ into a commutative ring $W_r(A)$, the ring of \textbf{$p$-typical
Witt vectors of length $r$ over $A$}, with zero $(0, \dots, 0)$ and unit
$(1, 0, \dots, 0)$. Additive inverses are likewise given by universal integral
polynomials; we write $\ominus$ for Witt subtraction. The \textbf{ghost map}
\[
w \colon W_r(A) \longrightarrow A^r,
\qquad
\vec a \longmapsto \bigl(w_0(\vec a), w_1(\vec a), \dots, w_{r-1}(\vec a)\bigr)
\]
is a ring homomorphism to the product ring $A^r$; it is injective when $p$ is a
non-zero-divisor in $A$, and bijective when $p \in A^\times$.
\end{theorem}

We index the components of a Witt vector $\vec a = (a_0, \dots, a_{r-1})$ by
$0, \dots, r-1$. For $r = 1$ the polynomials reduce to $S_0 = X_0 + Y_0$ and
$P_0 = X_0 Y_0$, so $W_1(A) = A$ as a ring.

Since the ring operations are given by universal polynomials with integer
coefficients, $W_r$ is a functor: a ring homomorphism $\sigma \colon A \to B$ induces
the ring homomorphism
\[
W_r(\sigma) \colon W_r(A) \to W_r(B),
\qquad
W_r(\sigma)(a_0, \dots, a_{r-1}) = (\sigma a_0, \dots, \sigma a_{r-1}),
\]
acting \emph{componentwise}. Two consequences of this trivial-looking observation are
used repeatedly and deserve emphasis: a ring automorphism $\sigma$ of $A$ commutes
with $\oplus$ and $\ominus$ on $W_r(A)$, and a Witt vector $\vec a \in W_r(A)$
satisfies $W_r(\sigma)(\vec a) = \vec a$ if and only if \emph{each component} $a_n$ is
$\sigma$-invariant. In particular, for a group $\Gamma$ acting on $A$ by ring
automorphisms,
\begin{equation}\label{eq:witt_invariants_componentwise}
W_r(A)^{\Gamma} \;=\; W_r\bigl(A^{\Gamma}\bigr)
\qquad \text{inside } W_r(A).
\end{equation}

The uniqueness statement in \Cref{theorem:witt_ring_structure} is an instance of a
general principle that we invoke several times, so we record it explicitly.

\begin{lemma}[ghost principle]\label{lemma:ghost_principle}
Let $\Phi$ and $\Psi$ be two operations on Witt vectors given by universal polynomials
with integer coefficients in the components of their arguments. If
$w \circ \Phi = w \circ \Psi$ holds over the polynomial ring
$\Z[X_0, X_1, \dots]$ on the components of the arguments, then $\Phi = \Psi$ over
every commutative ring.
\end{lemma}

\begin{proof}
Over $\Z[X_0, X_1, \dots]$ the prime $p$ is a non-zero-divisor, so the ghost map is
injective by \Cref{theorem:witt_ring_structure} and $\Phi = \Psi$ there; this is an
identity of integral polynomials, which persists under every base change
$\Z[X_0, X_1, \dots] \to A$.
\end{proof}

The only structural feature of the addition polynomials $S_n$ that we shall need is
their ``triangular'' shape:

\begin{prop}[shape of Witt addition]\label{prop:witt_addition_shape}
For every $n \geq 0$ there is a polynomial
\[
C_n \in \Z[X_0, \dots, X_{n-1}; Y_0, \dots, Y_{n-1}]
\]
involving only the variables of index \emph{strictly less} than $n$, such that
\begin{equation}\label{eq:witt_addition_shape}
(\vec a \oplus \vec b)_n \;=\; a_n + b_n + C_n(a_0, \dots, a_{n-1}; b_0, \dots, b_{n-1}).
\end{equation}
We refer to $C_n$ as the \textbf{carry polynomial}. The analogous statement holds for
$\ominus$.
\end{prop}

\begin{proof}
By induction on $n$, using the recursion defining $S_n$:
\[
p^n S_n \;=\; \sum_{j \leq n} p^j \bigl(X_j^{p^{n-j}} + Y_j^{p^{n-j}}\bigr)
\;-\; \sum_{j < n} p^j S_j^{p^{n-j}} .
\]
The terms with $j = n$ contribute $p^n (X_n + Y_n)$, and all remaining terms involve
only variables and polynomials $S_j$ of index $j < n$, which by induction lie in
$\Z[X_{<n}; Y_{<n}]$. Since $S_n$ has integer coefficients
(\Cref{theorem:witt_ring_structure}), $C_n := S_n - X_n - Y_n$ is an integral
polynomial in the variables of index $< n$. The statement for $\ominus$ follows by the
same argument applied to the inversion polynomials.
\end{proof}

\subsection{Frobenius, Verschiebung, and truncation}

From now on, $A$ denotes an $\F_p$-algebra; this is the only case we use. Three
operators act on the rings $W_r(A)$:

\begin{definition}\label{definition:witt_operators}
Let $A$ be an $\F_p$-algebra.
\begin{enumerate}
    \item The \textbf{Frobenius} $F \colon W_r(A) \to W_r(A)$ is
    $F := W_r(\varphi_A)$, where $\varphi_A \colon a \mapsto a^p$ is the absolute
    Frobenius of $A$. Explicitly, $F$ acts componentwise:
    $F(a_0, \dots, a_{r-1}) = (a_0^p, \dots, a_{r-1}^p)$. By functoriality, $F$ is a
    ring endomorphism of $W_r(A)$.
    \item The \textbf{Verschiebung} is the shift
    \[
    V \colon W_{r-1}(A) \to W_r(A),
    \qquad
    V(a_0, \dots, a_{r-2}) = (0, a_0, \dots, a_{r-2}).
    \]
    Iterating, $V^j \colon W_{r-j}(A) \to W_r(A)$; in particular
    $V^{r-1} \colon A = W_1(A) \to W_r(A)$ sends $c \mapsto (0, \dots, 0, c)$.
    \item The \textbf{truncation} is
    \[
    t \colon W_r(A) \to W_{r-1}(A),
    \qquad
    t(a_0, \dots, a_{r-1}) = (a_0, \dots, a_{r-2}).
    \]
\end{enumerate}
\end{definition}

\begin{rem*}
For a general (not necessarily $\F_p$-) algebra $A$, the Frobenius is defined by
universal polynomials characterized by $w_n \circ F = w_{n+1}$, and the componentwise
formula above holds precisely in characteristic $p$; see
\cite[II, \S6]{serreLF} or \cite[\S2]{Rabinoff2014Witt}. We will not need the general
case.
\end{rem*}

We record from the general theory (\cite[II, \S6]{serreLF},
\cite[\S2]{Rabinoff2014Witt}) how multiplication by $p$ interacts with these
operators.

\begin{prop}\label{prop:witt_multiplication_by_p}
Let $A$ be an $\F_p$-algebra. Then:
\begin{enumerate}
    \item $t$ is a natural ring homomorphism commuting with $F$, and its kernel is
    \[
    \ker\bigl(t \colon W_r(A) \to W_{r-1}(A)\bigr)
    \;=\; V^{r-1} A \;=\; \{(0, \dots, 0, c) : c \in A\}.
    \]
    \item $V$ is additive, and $F V = V F$.
    \item Multiplication by $p$ on $W_r(A)$ equals $V \circ F = F \circ V$
    (composed with the appropriate truncation); explicitly,
    \[
    p \cdot (a_0, \dots, a_{r-1}) \;=\; (0,\, a_0^p,\, \dots,\, a_{r-2}^p).
    \]
    Iterating, $p^j \cdot \vec a = (0, \dots, 0, a_0^{p^j}, \dots, a_{r-1-j}^{p^j})$;
    in particular $p^r = 0$ on $W_r(A)$, and
    \begin{equation}\label{eq:witt_p_to_r_minus_1}
    p^{r-1} \cdot (a_0, \dots, a_{r-1}) \;=\; \bigl(0, \dots, 0,\, a_0^{p^{r-1}}\bigr).
    \end{equation}
\end{enumerate}
\end{prop}

\begin{proof}
(1) That $t$ is a ring homomorphism follows from the recursive definition of the
$S_n, P_n$: the first $r-1$ components of a sum or product depend only on the first
$r-1$ components of the arguments (\Cref{prop:witt_addition_shape}). Naturality and
the commutation with the (componentwise) $F$ are clear, and the kernel is computed
componentwise. (2) Additivity of $V$ follows from \Cref{lemma:ghost_principle} and
the ghost identity $w_n(V\vec a) = p\, w_{n-1}(\vec a)$ (with $w_{-1} := 0$); the
commutation $FV = VF$ is immediate since $F$ is componentwise. (3) The identity
$F V = p$ holds over arbitrary rings for the general Frobenius, by
\Cref{lemma:ghost_principle} applied to
$w_n(FV\vec a) = w_{n+1}(V \vec a) = p\, w_n(\vec a) = w_n(p \vec a)$; in
characteristic $p$ the general Frobenius is componentwise, giving the displayed
formula. See \cite[II, \S6, Proposition~8 ff.]{serreLF} for details.
\end{proof}

\begin{example}\label{example:witt_of_Fp}
$W_r(\F_p) \cong \Z/p^r\Z$ canonically. Indeed, by
\Cref{eq:witt_p_to_r_minus_1} the unit $\underline 1 = (1, 0, \dots, 0)$ satisfies
$p^j \cdot \underline 1 = (0, \dots, 0, 1, 0, \dots, 0)$ ($1$ in position $j$), so
$\underline 1$ has additive order $p^r$; since $|W_r(\F_p)| = p^r$, the additive group
is cyclic of order $p^r$ generated by the unit, and $m \mapsto m \cdot \underline 1$
is a ring isomorphism $\Z/p^r\Z \xrightarrow{\;\cong\;} W_r(\F_p)$. We write
$\underline m \in W_r(\F_p)$ for the image of $m \in \Z/p^r\Z$. (Similarly, the full
ring of Witt vectors satisfies $W(\F_p) \cong \Z_p$; we shall not need this.)
\end{example}

\section{Artin--Schreier--Witt theory}
\label{appendix:asw_development}

We retain the notation for Witt vectors from the beginning of this appendix.
For an $\F_p$-algebra $A$, Frobenius acts componentwise on $W_r(A)$.

\subsection{The ASW operator and bookkeeping}

\begin{definition}\label{definition:asw_operator}
The \textbf{Artin--Schreier--Witt operator} is
\[
\wp:=F-\mathrm{id},\qquad
\wp\vec u=F\vec u\ominus\vec u.
\]
It is an endomorphism of the additive group of $W_r(A)$ and commutes with
truncation and Verschiebung.
\end{definition}

\begin{lemma}[Witt bookkeeping]\label{lemma:witt_bookkeeping}
Let $A$ be an $\F_p$-algebra.
\begin{enumerate}[label=(\roman*)]
    \item The truncation $t:W_r(A)\to W_{r-1}(A)$ is a ring homomorphism
    commuting with $F$, hence with $\wp$, and
    $\ker(t)=V^{r-1}A$.
    \item The map $V^{r-1}:A\to W_r(A)$ is additive and
    $\wp\circ V^{r-1}=V^{r-1}\circ\wp$.
    \item For every $\vec a\in W_r(A)$ and $c\in A$,
    \[
    \vec a\oplus V^{r-1}(c)
    =(a_0,\dots,a_{r-2},a_{r-1}+c).
    \]
\end{enumerate}
\end{lemma}

\begin{proof}
The first two assertions follow from
\Cref{prop:witt_multiplication_by_p}. For (iii), consider the two operations
\[
\begin{aligned}
 \Phi(\vec X,C)
 &:=\vec X\oplus V^{r-1}(C)
   =\bigl(S_n(\vec X;0,\dots,0,C)\bigr)_{0\leq n\leq r-1},\\
 \Psi(\vec X,C)
 &:=(X_0,\dots,X_{r-2},X_{r-1}+C).
\end{aligned}
\]
Both are given by universal polynomials with integer coefficients, as required
in \Cref{lemma:ghost_principle}. Over
$\Z[X_0,\dots,X_{r-1},C]$, for $n<r-1$ we have
\[
 w_n\bigl(\Phi(\vec X,C)\bigr)
 =w_n(\vec X)+w_n\bigl(V^{r-1}(C)\bigr)
 =w_n(\vec X)
 =w_n\bigl(\Psi(\vec X,C)\bigr).
\]
For the final ghost component,
\[
\begin{aligned}
 w_{r-1}\bigl(\Phi(\vec X,C)\bigr)
 &=w_{r-1}(\vec X)+p^{r-1}C\\
 &=\sum_{i=0}^{r-2}p^iX_i^{p^{r-1-i}}
   +p^{r-1}(X_{r-1}+C)\\
 &=w_{r-1}\bigl(\Psi(\vec X,C)\bigr).
\end{aligned}
\]
Thus $w\circ\Phi=w\circ\Psi$, so
\Cref{lemma:ghost_principle} gives $\Phi=\Psi$ over every commutative ring,
and in particular over $A$.
\end{proof}

The last assertion is special to the top Verschiebung slice. Witt addition is
not componentwise in lower positions, and the carry polynomials in
\Cref{eq:witt_addition_shape} cannot in general be discarded.

\subsection{Torsors and classification}
The Witt addition polynomials make $\A^r_{\F_p}$ into the additive Witt
vector group scheme $\bW_r$. We distinguish this group scheme from its
groups of points: for an $\F_p$-algebra $A$, its group of $A$-valued points
is $\bW_r(A)=W_r(A)$. If $X$ is an $\F_p$-scheme, we write
$W_r(\F_p)_X$ for the constant group scheme over $X$ associated with the
abstract group $W_r(\F_p)$. The operator $\wp=F-\mathrm{id}$ is an
endomorphism of $\bW_r$ and fits into an exact sequence of group schemes for
the \'etale topology
\begin{equation}\label{eq:asw_lang_sequence}
0\longrightarrow W_r(\F_p)_{\F_p}\cong(\Z/p^r\Z)_{\F_p}
\longrightarrow\bW_r\xrightarrow{\wp}\bW_r\longrightarrow0.
\end{equation}

The following is the Witt-vector analogue of
\Cref{lemma:classical_cyclic_equations_etale_torsors}. Its coordinate proof
requires the triangular form of Witt addition because, unlike the two
classical equations, Witt subtraction is not componentwise.

\begin{lemma}[ASW covers are \'etale torsors]
\label{lemma:asw_etale_torsor}
Let $A$ be an $\F_p$-algebra and $\vec g\in W_r(A)$. Let
$\vec X=(X_0,\dots,X_{r-1})$ denote the universal Witt vector of length $r$
over $A$, so that $A[\vec X]=A[X_0,\dots,X_{r-1}]$. The fiber of $\wp$
over $\vec g$ is
\[
T_{\vec g}=\Spec B,
\qquad
B:=A[\vec X]\big/\bigl(\wp\vec X\ominus\vec g\bigr),
\]
where the parenthesized Witt vector denotes the ideal generated by its $r$
coordinate polynomials.
Here $\vec g\in W_r(A)=\bW_r(A)$ is viewed as a morphism
$\Spec A\to\bW_r$, and $T_{\vec g}$ is defined by the Cartesian square
\[
\begin{tikzcd}
T_{\vec g} \arrow[r] \arrow[d] & \bW_r \arrow[d,"\wp"] \\
\Spec A \arrow[r,"\vec g"'] \arrow[dr] & \bW_r \arrow[d] \\
& \Spec\F_p,
\end{tikzcd}
\]
where the maps to $\Spec\F_p$ are the structure maps.
The algebra $B$ is finite free over $A$, with basis
$\{\prod_nX_n^{a_n}:0\leq a_n\leq p-1\}$, and is \'etale over $A$.
Translation by $W_r(\F_p)\cong\Z/p^r\Z$, explicitly
$\vec X\mapsto\vec X\oplus\vec c$, makes $T_{\vec g}$ a
$\Z/p^r\Z$-torsor over $\Spec A$.
\end{lemma}

\begin{proof}
By the triangular shape of Witt addition
(\Cref{eq:witt_addition_shape}), the $n$-th equation in
$\wp\vec X\ominus\vec g=0$ has the form
\[
X_n^p-X_n-g_n+C_n'(X_0,\dots,X_{n-1};g_0,\dots,g_{n-1})=0.
\]
Starting with $A$, first adjoin $X_0$ subject to the equation of index $0$.
Then, for each $1\leq n<r$, adjoin $X_n$ subject to the displayed equation
of index $n$; its coefficients lie in the algebra obtained at the preceding
step. Because this equation is monic of degree $p$ in $X_n$, the $n$-th step
is free of rank $p$, with basis $1,X_n,\dots,X_n^{p-1}$. After all $r$ steps
we obtain $B$, finite free over $A$ with the stated basis. The Jacobian matrix
is lower triangular with diagonal entries $-1$, so $B/A$ is \'etale.
Moreover, this cover is the
pullback of $\wp:\bW_r\to\bW_r$ along the section $\vec g$. The kernel of
the group-scheme morphism $\wp$ is the constant group scheme
$W_r(\F_p)_{\F_p}$, characterized by $F\vec u=\vec u$. Translation by the
base change of this kernel to $\Spec A$ preserves the fiber because, for
$\vec c\in\ker(\wp)$,
\[
\wp(\vec X\oplus\vec c)=\wp(\vec X)\oplus\wp(\vec c)=\wp(\vec X).
\]
Any two points in the same fiber differ by a unique such $\vec c$, so this
translation action makes the pullback a torsor.
\end{proof}

\begin{theorem}[Artin--Schreier--Witt theory;
{\cite[\S3--4]{Thomas2005}}]
\label{theorem:asw_theory}
Let $M$ be a field of characteristic $p$, and let $r\geq1$. For
$\vec g\in W_r(M)$, let
\[
T_{\vec g}:=
\Spec\!\left(M[\vec X]\big/\bigl(\wp\vec X\ominus\vec g\bigr)\right)
\]
be the $\Z/p^r\Z$-torsor of \Cref{lemma:asw_etale_torsor}. The assignment
$[\vec g]\mapsto[T_{\vec g}]$ is a canonical isomorphism of abelian groups
\[
W_r(M)/\wp W_r(M)\;\xrightarrow{\ \cong\ }\;
H^1_{\mathrm{\acute et}}
\bigl(\Spec M,(\Z/p^r\Z)_{\Spec M}\bigr).
\]
Thus this quotient classifies $\Z/p^r\Z$-torsors over $\Spec M$.

If $[\vec g]$ has order $p^s$, where $0\leq s\leq r$, then $T_{\vec g}$
has $p^{r-s}$ connected components, each noncanonically isomorphic to
$\Spec L$ for the same cyclic extension $L/M$ of degree $p^s$.
Let $H\subseteq\Z/p^r\Z$ be the unique subgroup of order $p^s$.
Each component is an $H$-torsor, and $T_{\vec g}$ is induced from any one
of them. In particular, $T_{\vec g}$ is connected if and only if
$[\vec g]$ has order $p^r$. Every cyclic extension of $M$ of degree
dividing $p^r$ occurs in this way.
\end{theorem}

\begin{proof}
Let $M^s$ be a separable closure of $M$, and put $\Gamma=\Gal(M^s/M)$.
The triangular equations in the proof of
\Cref{lemma:asw_etale_torsor} show that $\wp$ is surjective on $W_r(M^s)$
and has kernel $W_r(\F_p)\cong\Z/p^r\Z$. Choose $\vec u\in W_r(M^s)$
with $\wp\vec u=\vec g$. Under the torsor--character correspondence
of \Cref{cor:abelian_equivilance_fundamental_torsors}, $T_{\vec g}$
corresponds to
\[
\chi_{\vec g}\colon\Gamma\longrightarrow W_r(\F_p),
\qquad
\chi_{\vec g}(\sigma)=\sigma(\vec u)\ominus\vec u.
\]
This is the character of the connecting class $[T_{\vec g}]$ in
\Cref{eq:common_torsor_connecting_map}.

We first check that this map is well defined. Any other solution of
$\wp\vec u=\vec g$ is $\vec u\oplus\vec c$ for some
$\vec c\in W_r(\F_p)$. Since $\Gamma$ fixes $\vec c$, this change does not
alter $\chi_{\vec g}$. Similarly, replacing $\vec g$ by
$\vec g\oplus\wp\vec v$, with $\vec v\in W_r(M)$, replaces $\vec u$ by
$\vec u\oplus\vec v$ and again leaves the character unchanged. If
$\chi_{\vec g}=0$, then $\vec u$ is fixed by $\Gamma$, so
$\vec u\in W_r(M)$ by \Cref{eq:witt_invariants_componentwise}. Hence
$[\vec g]=0$, and the map is injective.

It remains to prove surjectivity. We claim that
\[
H^1_{\mathrm{cont}}\bigl(\Gamma,W_r(M^s)\bigr)=0.
\]
For $r=1$, this is additive Hilbert~90. For $r>1$,
\Cref{lemma:witt_bookkeeping}(i) gives a short exact sequence of
$\Gamma$-modules
\[
0\longrightarrow M^s\xrightarrow{V^{r-1}}W_r(M^s)
\xrightarrow{t}W_{r-1}(M^s)\longrightarrow0.
\]
The associated long exact sequence contains
\[
H^1_{\mathrm{cont}}(\Gamma,M^s)
\longrightarrow
H^1_{\mathrm{cont}}\bigl(\Gamma,W_r(M^s)\bigr)
\longrightarrow
H^1_{\mathrm{cont}}\bigl(\Gamma,W_{r-1}(M^s)\bigr).
\]
The outer groups vanish by additive Hilbert~90 and induction, so the middle
group vanishes. The long exact sequence associated with
\Cref{eq:asw_lang_sequence} now shows that the connecting map is surjective.
This proves the claimed isomorphism.

Now suppose $[\vec g]$ has order $p^s$, and put
$H=\operatorname{im}(\chi_{\vec g})$. The isomorphism preserves orders,
so $\chi_{\vec g}$ has order $p^s$. Since $H$ is cyclic, the order of the
character equals $|H|$. By \Cref{theorem:equivalence_fundamental_covers},
the connected components of $T_{\vec g}$ correspond to the $\Gamma$-orbits
on its geometric fiber. Choosing $\vec u$ identifies that fiber with
$W_r(\F_p)$, with $\Gamma$ acting by translation through $\chi_{\vec g}$.
The orbits are the cosets of $H$, so there are $p^{r-s}$ components.
Each geometric point has stabilizer $\ker(\chi_{\vec g})$ in $\Gamma$.
Thus each component is isomorphic to $\Spec L$, where
$L=(M^s)^{\ker(\chi_{\vec g})}$ is cyclic of degree $|H|=p^s$ over $M$.
The subgroup $H$ preserves each component and acts simply transitively on
its geometric fiber. Each component is therefore an $H$-torsor, and
inducing along $H\hookrightarrow W_r(\F_p)$ recovers $T_{\vec g}$.

Conversely, let $L/M$ be cyclic of degree $p^s$ dividing $p^r$, and embed
$L$ in $M^s$. Choose an injection $\Gal(L/M)\hookrightarrow W_r(\F_p)$
and compose it with the
restriction map $\Gamma\to\Gal(L/M)$. By the surjectivity proved above, the
resulting character is $\chi_{\vec g}$ for some $\vec g\in W_r(M)$. The
fixed field of its kernel is $L$, so the connected components of $T_{\vec g}$ are
isomorphic to $\Spec L$.
\end{proof}

% [PROPOSED REWRITE, for comparison with the theorem and proof above]
{\color{green!50!black}
\begin{theorem*}[Proposed rewrite of \Cref{theorem:asw_theory}]
Let $M$ be a field of characteristic $p$, and let $r\geq1$. For
$\vec g\in W_r(M)$ put
$T_{\vec g}:=\Spec M[\vec X]/(\wp\vec X\ominus\vec g)$.
\begin{enumerate}[label=(\roman*)]
  \item The map $[\vec g]\mapsto[T_{\vec g}]$ is a group isomorphism
  \[
  W_r(M)/\wp W_r(M)\xrightarrow{\ \cong\ }
  H^1_{\mathrm{\acute et}}\bigl(\Spec M,(\Z/p^r\Z)_{\Spec M}\bigr).
  \]
  The target is \'etale cohomology with coefficients in the constant group
  scheme $(\Z/p^r\Z)_{\Spec M}\cong W_r(\F_p)_{\Spec M}$
  (\Cref{example:witt_of_Fp}). Its elements are the classes of
  $\Z/p^r\Z$-torsors over $\Spec M$ (see the note in
  \Cref{section:torsors_and_cohomology_note}).
  \item Suppose $[\vec g]$ has order $p^s$, and let $H\subseteq\Z/p^r\Z$
  be the unique subgroup of order $p^s$. Then $T_{\vec g}$ is a disjoint
  union of $p^{r-s}$ copies of $\Spec L$, where $L/M$ is cyclic of degree
  $p^s$. Each copy is an $H$-torsor, and $T_{\vec g}$ is induced from it.
  In particular, $T_{\vec g}$ is connected if and only if $[\vec g]$ has
  order $p^r$.
  \item Every cyclic extension of $M$ of degree dividing $p^r$ is such
  an $L$.
\end{enumerate}
\end{theorem*}

\begin{proof}[Proposed proof]
\emph{(i), injectivity.} Apply \Cref{lemma:asw_etale_torsor} to the
universal vector over $\F_p[g_0,\dots,g_{r-1}]$: the map
$\wp:\bW_r\to\bW_r$ is finite \'etale and surjective, with kernel
$W_r(\F_p)\cong\Z/p^r\Z$. So \Cref{eq:common_torsor_connecting_map}
gives an injective homomorphism $[\vec g]\mapsto[T_{\vec g}]$.

\emph{(i), surjectivity.} The long exact sequence in
\Cref{appendix:common_group_scheme_construction} continues with
$H^1(M,\Z/p^r\Z)\to H^1_{\mathrm{\acute et}}(\Spec M,\bW_r)$. So it is
enough to show that $H^1_{\mathrm{\acute et}}(\Spec M,\bW_r)=0$. For $r=1$
this is additive Hilbert~90 \cite[X, \S1]{serreLF}. For $r>1$, the sequence
$0\to\bG_a\xrightarrow{V^{r-1}}\bW_r\xrightarrow{t}\bW_{r-1}\to0$
(\Cref{prop:witt_multiplication_by_p}) is exact. Its outer terms have zero
$H^1$ by induction, so the middle one does too.

\emph{(ii).} Let $\chi:\Gal(M^s/M)\to\Z/p^r\Z$ be the character of
$T_{\vec g}$ (\Cref{cor:abelian_equivilance_fundamental_torsors}). The map
in (i) is a group isomorphism, so $\chi$ has order $p^s$. A subgroup of
$\Z/p^r\Z$ is cyclic, so the order of $\chi$ equals the size of its image.
Hence $\operatorname{im}\chi=H$. By
\Cref{theorem:equivalence_fundamental_covers}, $T_{\vec g}$ is induced from
the connected $H$-torsor $\Spec L$, where $L$ is the fixed field of
$\ker\chi$ and $[L:M]=|H|=p^s$. An induced torsor has
$[\Z/p^r\Z:H]=p^{r-s}$ components.

\emph{(iii).} Let $L/M$ be cyclic of degree $p^s\mid p^r$. Choose an
injection $\Gal(L/M)\hookrightarrow\Z/p^r\Z$. The composite
$\Gal(M^s/M)\to\Z/p^r\Z$ is a character. By (i) it comes from some
$\vec g$, and by (ii) the components of $T_{\vec g}$ are $\Spec L$.
\end{proof}
}

\begin{cor}\label{cor:asw_class_order_and_first_component}
Let $M\supseteq\F_p$ and let $\vec g\in W_r(M)$ have $g_0=0$. Then
\[
p^{r-1}[\vec g]=0 \qquad\text{in } W_r(M)/\wp W_r(M).
\]
Consequently, a class of
order $p^r$ has nonzero zeroth component in every representative.
\end{cor}

\begin{proof}
By \Cref{eq:witt_p_to_r_minus_1},
$p^{r-1}\vec g=(0,\dots,0,g_0^{p^{r-1}})=0$ already in $W_r(M)$.
\end{proof}

\subsection{Reduced polynomial representatives}

\begin{definition}\label{definition:reduced_polynomial_vector}
A vector $\vec g=(g_0,\dots,g_{r-1})\in W_r(k((x)))$ is a
\textbf{reduced polynomial vector} if it is reduced and every
$g_j\in x^{-1}k[x^{-1}]$.
\end{definition}

\begin{lemma}[Surjectivity on integral Witt vectors]
\label{lemma:wp_surjective_on_integral_witt}
If $k$ is algebraically closed, then
\[
\wp:W_r(k[[x]])\longrightarrow W_r(k[[x]])
\]
is surjective.
\end{lemma}

\begin{proof}
Proceed by induction on $r$. For $r=1$, given
$h=\sum_{i\geq0}b_ix^i$, choose $u_0$ with
$u_0^p-u_0=b_0$ and define the remaining coefficients recursively by
$u_i=u_{i/p}^p-b_i$, with $u_{i/p}=0$ when $p\nmid i$. Then
$u^p-u=h$.

For $r>1$, choose by induction $\vec u'\in W_{r-1}(k[[x]])$ with
$\wp\vec u'=t(\vec h)$ and lift it to
$\widetilde u=(\vec u',0)$. By \Cref{lemma:witt_bookkeeping}(i),
$\wp\widetilde u\ominus\vec h=V^{r-1}(g)$ for some $g\in k[[x]]$.
Choose $w\in k[[x]]$ with $\wp w=-g$. Then
\Cref{lemma:witt_bookkeeping}(ii) gives
$\wp(\widetilde u\oplus V^{r-1}w)=\vec h$.
\end{proof}

\begin{prop}[Reduced polynomial representatives]
\label{prop:reduced_polynomial_representatives}
Every class in $W_r(k((x)))/\wp W_r(k((x)))$ has a reduced polynomial
representative.
\end{prop}

\begin{proof}
Proceed by induction on $r$. For $r=1$, write a Laurent series as the sum of
its principal part and an element of $k[[x]]$. The latter can be removed by
\Cref{lemma:wp_surjective_on_integral_witt}. If the remaining principal part
has leading term $cx^{-m}$ with $p\mid m$, subtract
\[
\wp(c^{1/p}x^{-m/p})=cx^{-m}-c^{1/p}x^{-m/p}.
\]
This strictly decreases the pole order; iteration produces zero or a pole
order prime to $p$.

For the inductive step, first reduce the truncation of the given vector and
lift the change of representative to length $r$. The first $r-1$ components
are then reduced polynomials. By
\Cref{lemma:witt_bookkeeping}(iii), subtracting an element of the form
$\wp(V^{r-1}z)=V^{r-1}(\wp z)$ changes only the last component. Apply the
length-one procedure to that component: remove its integral part and then
successively remove leading terms whose pole orders are divisible by $p$.
This leaves all earlier components unchanged and produces a reduced polynomial
vector representing the original class.
\end{proof}


\section{Weighted Witt addition in the cyclic calculation}

The following homogeneity property of the universal addition polynomials is what
makes degree estimates on Witt sums possible: even though the carry polynomials $C_n$
of \Cref{eq:witt_addition_shape} are complicated, they are \emph{isobaric} for the
weighting in which the $j$-th component has weight $p^j$.

\begin{lemma}[isobaricity of Witt addition]\label{lemma:witt_addition_isobaric}
Assign to the variables $X_j$ and $Y_j$ the weight $p^j$. Then the $n$-th addition
polynomial $S_n$ is isobaric of weight $p^n$: every monomial of $S_n$ has total weight
exactly $p^n$. The same holds for each component of a $d$-fold iterated Witt sum
$\bigoplus_{i=1}^{d} \vec a_i$, viewed as a universal polynomial in the components
$a_{i,j}$ of weight $p^j$.
\end{lemma}

\begin{proof}
The $n$-th ghost polynomial $w_n = \sum_{j \le n} p^j X_j^{p^{n-j}}$ is isobaric of
weight $p^n$ for this weighting. The recursion
\[
p^n S_n = \sum_{j \le n} p^j \bigl(X_j^{p^{n-j}} + Y_j^{p^{n-j}}\bigr)
- \sum_{j < n} p^j S_j^{p^{n-j}}
\]
preserves isobaricity of weight $p^n$ by induction on $n$: each $S_j$ is isobaric of
weight $p^j$, so $S_j^{p^{n-j}}$ is isobaric of weight $p^n$. Isobaricity persists
under iterated substitution of Witt sums into Witt sums (substituting an isobaric
polynomial of weight $p^j$ for a variable of weight $p^j$ preserves weights), giving
the statement for $d$-fold sums.
\end{proof}
